% TODO:
% -

\documentclass[twocolumn]{AASTeX7/aastex701}

\shorttitle{\texttt{harv}: fast companion orbit inference}
\shortauthors{Price-Whelan et al.}

\usepackage{preamble/preamble-adrn}

\newcommand{\harv}{\package{harv}}
\newcommand{\joker}{\textsl{The Joker}}

\begin{document}

\title{
  \texttt{harv}: \\
  Fast companion orbit inference with radial velocities and \gaia\ epoch astrometry
}

\author[orcid=0000-0003-0872-7098]{Adrian M. Price-Whelan}
\email{aprice-whelan@flatironinstitute.org}
\affiliation{
  Center for Computational Astrophysics, Flatiron Institute,
  162 Fifth Ave, New York, NY 10010, USA
}

\begin{abstract}
Population studies of stellar companions (binary stars, exoplanets, compact objects)
often start from a sample of systems with well-constrained orbits.
Inferring population statistics from such a sample requires modeling the detection
efficiency of the orbit-fitting pipeline along with the survey selection function,
usually through injection/recovery tests.
An alternative approach is to include every source from the parent sample in the
population inference.
This requires generating orbital parameter posterior samples even for sources with
sparse, noisy, or null signals in the same manner as the well-constrained systems.
\harv\ is an open-source Python package that enables posterior sampling with radial
velocities (RVs), \gaia\ epoch astrometry, or both, across a range of data quality
regimes.
Like its predecessor \joker, \harv\ enables posterior sampling of Keplerian orbit
parameters either through traditional Markov chain Monte Carlo (MCMC) or via a custom
rejection sampler that analytically marginalizes over linear parameters in the model to
improve efficiency.
For sparse or uninformative data, the posterior is often highly multimodal, which makes
MCMC inefficient and unreliable, but the rejection sampler still returns independent
posterior samples for such cases.
\harv\ is built on \jax, which provides just-in-time compilation, CPU and GPU
execution, and automatic differentiation for gradient-based inference methods.
We demonstrate \harv\ on (1) sparse \apogee\ RVs with multimodal posteriors, (2) the
planet Gaia-4b, jointly fitting RVs and \gaia\ \dr{4} pre-release astrometry, and (3)
the active planet host HD~285968, using a Gaussian process to model the stellar
variability.
\harv\ is designed to both be flexible, and to scale to the millions of sources with
time-domain data from \gaia\ \dr{4}, \sdssv, and future stellar surveys.
\end{abstract}

% =====================================================================================

%--------------------------------------------------------------------
\section{Introduction}
\label{sec:intro}

Stars, most often, have at least one companion --- another star, a brown dwarf, planet,
or compact object, for example \citep[e.g.,][]{TODO, Moe}.
The distribution of companion properties (e.g., mass, orbital period, eccentricity,
etc.), their occurrence rates, and their dependence on the properties of the host star
all provide important constraints on models of star formation, binary interactions,
planet formation, and stellar evolution \citep{TODO}.
These systems are also important benchmarks for stellar astrophysics \citep{TODO}.
Measuring distributions of companion properties requires time-domain survey data and
inferences of the orbital parameters of large samples of systems.

Population studies of stellar companions often start by selecting systems with
well-constrained orbits from larger catalogs.
For example, a typical pipeline will fit an orbit to every source in some survey sample,
keep the orbital solutions that pass some cuts on significance or goodness of fit, and
then use the resulting sample to infer population statistics.
An example of this approach is the way non-single stars (NSS) are processed in the
\gaia\ \dr{3} astrometric binary catalog, in which a decision is made for each source
whether to fit a single-star or binary-star model, and only stars that pass a
significance cut on the binary model are included in the final catalog
\citep{Halbwachs:2023, GaiaDR3Multiplicity}.
However, significance cuts often depend on the orbital parameters themselves, so to turn
a sample defined in this way into a robust population inference, one must  model both
the parent catalog selection function \textsl{and} the detection efficiency of the
orbit-fitting-and-selection pipeline, which often requires expensive injection and
recovery tests \citep[e.g.,][]{Cumming:2008, Fulton:2021, ElBadry:2024}.
\todo{check these and maybe add references}

An alternate approach to stellar companion population inference is to instead include
all sources from the parent catalog in the model, whether or not their orbits are
precisely constrained.
In the context of a hierarchical Bayesian model (hereafter ``hierarchical model''), the
likelihood function $\mathcal{L}$ for the population parameters $\vec{\alpha}$ (e.g.,
parameters that describe the companion mass function, period distribution, etc.) follows
a hierarchical structure in which the full likelihood is a product over the individual
likelihoods for each source $n$ with orbital parameters $\vec{\theta}_n$ in the parent
catalog:
\begin{align}
  \mathcal{L}(\vec{\alpha} \,;\, \{D_n\}, \{\vec{\theta}_n\}) &=
    p(\{D_n\}, \{\vec{\theta}_n\} \given \vec{\alpha}) \\
    &= \prod_n p(D_n \given \vec{\theta}_n) \, p(\vec{\theta}_n \given \vec{\alpha})
    \quad .
\end{align}
Here the per-system orbital parameters $\vec{\theta}_n$ are explicitly included in the
model, which means the dimensionality of the parameter space grows with the number of
sources in the parent catalog (and this is generally bad).
A common approach to make the hierarchical model more tractable is to marginalize over
the per-system parameters $\vec{\theta}_n$ to obtain a marginalized likelihood function
\begin{align}
  \mathcal{L}(\vec{\alpha} \,;\, \{D_n\}) &=
    p(\{D_n\} \given \vec{\alpha}) \\
  &= \prod_n
    \int \dd\vec{\theta}_n \,
      p(D_n \given \vec{\theta}_n) \, p(\vec{\theta}_n \given \vec{\alpha})
    \quad ,
\end{align}
but this integral is often not analytically solvable.
One can instead approximate the integral with Monte Carlo integration, which can be done
by drawing posterior samples of the per-system parameters $\vec{\theta}_n$ for each
source independently under an ``interim'' prior $p(\vec{\theta}_n \given
\vec{\alpha}_0)$, and then re-weighting the samples to the population prior
$p(\vec{\theta}_n \given \vec{\alpha})$ \citep{Hogg:2010, Foreman-Mackey:2014,
Price-Whelan:2020}:
\begin{equation}
  p(\{D_n\} \given \vec{\alpha}) \approx
    \prod_n \frac{Z_{{\rm int},n}}{K_n}
      \sum_{k=1}^{K_n}
        \frac{p(\vec{\theta}_{nk} \given \vec{\alpha})}
        {p(\vec{\theta}_{nk} \given \vec{\alpha}_0)}
  \quad ,
  \label{eq:hierarchical}
\end{equation}
where $Z_{{\rm int},n}$ is the evidence for the data $D_n$ of source $n$ under the
interim prior (which often can be ignored if the same interim prior is used for all
sources).

In the product in Equation~\ref{eq:hierarchical}, a source with uninformative data would
contribute samples that follow the interim prior and therefore do not constrain
$\vec{\alpha}$ (or do so weakly), but importantly the source still contributes to the
product.
In this framework, we do not need to make a decision about whether a source passes some
detection threshold, and therefore we do not need to model the detection efficiency of
our orbit fitting method (and therefore we do not need to do injection and recovery
tests).
The challenge (and computational cost) here instead is that we much generate posterior
samples of the orbital parameters $\vec{\theta}_n$ for every source in the parent
catalog, including sources with sparse or noisy data, or apparently null signals.
This is a challenge for Keplerian orbit fitting: with only a few radial velocity
measurements, or astrometry at low signal-to-noise, the per-source likelihood as a
function of period can have many narrow modes separated by wide regions of low
probability \citep[e.g.,][]{Price-Whelan:2017}.
Markov chain Monte Carlo (MCMC) methods struggle in this regime.

We previously developed \joker\ \citep{Price-Whelan:2017} for this regime with RV data
from the Apache Point Observatory Galactic Evolution Experiment (\apogee) in mind, which
often has only a few (but precisely-measured) RV measurements for hundreds of thousands
of stars.
\joker\ is a custom rejection sampler that can generate independent posterior samples of
Keplerian orbital parameters for a single source even when the posterior is highly
multimodal.
The trick we use to make rejection sampling tractable is that we analytically
marginalize over the linear parameters in the Keplerian model (the velocity
semi-amplitude, systemic velocity, and sometimes extra polynomial trend parameters) to
reduce the dimensionality of the parameter space that we need to sample over.
An advantage to rejection sampling is that the returned samples are independent draws
from the posterior, however many modes there are, and we do not need to worry about
convergence or autocorrelation analysis of chains (as we do with MCMC).
\joker\ has been used in a number of population studies of stellar companions
\citep[e.g.,][]{TODO}.

Our anticipation of two contemporary stellar surveys --- \gaia\ \dr{4} and the Milky Way
Mapper (MWM) program of \sdssv --- has motivated a ground-up redesign and
re-implementation of the core ideas of \joker.
\gaia\ \dr{4} \citep{TODO}, which will be released on 2026-12-02, will include a wealth
of time-series data, including epoch astrometry and radial velocities for tens of
millions to billions of stars.
MWM \citep{TODO} is continuing the \apogee\ survey legacy and will provide precise
multi-epoch RVs for millions of stars, most of which will have only a few visits.
For many of these sources, we will have both RVs and epoch astrometry, and to make the
most of the full samples we will need to be able to model the faint and therefore low
signal-to-noise companions.
Also, the data volumes expected from these surveys are immense, and we will need to be
able to generate posterior samples efficiently for millions of sources, which is
prohibitive for the current \joker\ implementation.

Here we present \harv: a successor to \joker\ and full re-implementation built on \jax\
\citep{Bradburg:2018} and \package{numpyro} \citep{Phan:2019}.
\harv\ is faster and more flexible than the original \joker\ implementation, and works
with both RVs and \gaia\ epoch astrometry (with support for more data types to come).
Because it is built on \jax, the same model code can be compiled and run on CPUs or
GPUs, and is differentiable, so can be used in gradient-based MCMC samplers or
variational inference methods instead of the built-in rejection sampler, when
appropriate.

Many packages and methods have been developed to efficiently sample or fit Keplerian
orbits to data.
For example, for RV data (\texttt{RadVel} \citealt{Fulton:2018}; \texttt{juliet} \citealt{Espinoza:2019}), for astrometry and combinations of data (\texttt{orvara} \citealt{Brandt:2021}; \texttt{orbitize!} \citealt{Blunt:2020}; \texttt{Octofitter} \citealt{Thompson:2023}; \texttt{kima} \citealt{Faria:2018}; \texttt{Exo-Striker} \citealt{Trifonov:2019}; \texttt{exoplanet} \citealt{Foreman-Mackey:2021}), and for forward modeling \gaia\ epoch astrometry and the \dr{3} orbit-fitting pipeline (\texttt{gaiamock}, \citealt{El-Badry:2024}).
These are designed mainly to characterize individual systems in which a companion has
been confidently detected, using MCMC or nested sampling.
The goal of \harv\ is to enable population inference with large samples of sources.

\todo{Mention below the Orvara marginalizations and that we use the kepler solver}


%--------------------------------------------------------------------
\section{Software design}
\label{sec:design}

\todo{Write this.}


%--------------------------------------------------------------------
\section{Inference}
\label{sec:inference}

\todo{Write this.}

%--------------------------------------------------------------------
\subsection{Marginalizing the linear parameters}
\label{sec:inference-linear}

%--------------------------------------------------------------------
\subsection{Rejection sampling}
\label{sec:inference-rejection}

%--------------------------------------------------------------------
\section{Case studies}
\label{sec:case-studies}

\todo{Write this.}

%--------------------------------------------------------------------
\subsection{Sparse APOGEE radial velocities}
\label{sec:cases-apogee}

\todo{Write this.}

%--------------------------------------------------------------------
\subsection{Gaia-4: joint radial velocities and DR4 astrometry}
\label{sec:cases-gaia4}

\todo{Write this.}

%--------------------------------------------------------------------
\subsection{HD 285968: a planet under stellar activity}
\label{sec:cases-hd285968}

\todo{Write this.}

%--------------------------------------------------------------------
\section{Discussion}
\label{sec:discussion}

\todo{Write this.}


%--------------------------------------------------------------------
\section{Summary}
\label{sec:summary}

\todo{Write this.}

\begin{acknowledgments}

  \todo{Write this.}

  Daniel Y, Sheila, Victor, Hogg, Foreman-Mackey, Will Thompson.

  \todo{LLM acknowledgement}

\begin{itemize}
  \item Flatiron Institute and CCA.
  \item Data acknowledgments: SDSS (\apogee), \gaia\ and DPAC, ESO and HARPS (programme
    IDs), and HPF, NEID and FIES as their data policies require.
\end{itemize}
\end{acknowledgments}

% \facilities{Gaia, Sloan (APOGEE), ESO:3.6m (HARPS), HET (HPF), WIYN (NEID), NOT (FIES)}

% TODO: software
% \software{harv, JAX, numpyro, equinox, unxt, quax, tinygp,
%           astropy, arviz, h5py, matplotlib, numpy, scipy, thejoker
%           }

\bibliographystyle{AASTeX7/aasjournalv7}
\bibliography{refs}{}

\end{document}
